\section{推理轨迹为何看起来像人：涌现、搜索与可验证边界}

有了中间 token 和 outcome-guided training，模型可能产生长而有结构的解题过程。课堂称之为 LLM reasoning 的“beauty”：不必显式编写穷举搜索器，token-to-token generation 也能出现分解、估计、回退和局部验证。这个现象值得研究，但最危险的误读是把语言表面直接当作认知同一性。

\subsection{从 token 生成中涌现出的结构}

slide 27 把现代 LLM 与经典人工智能的 exhaustive search 对比。更准确的理解是：模型把大量模式压进参数，并在生成时沿一条或少数轨迹展开；它未必显式枚举完整状态树。优势是推理路径紧凑、可与自然语言和工具接口结合；代价是覆盖不完备，且错误分支不一定能被系统性发现。

\lecturefigure{slide-27.jpg}{LLM reasoning 的吸引力：结构化过程可从 token 生成中出现}{27}

\teachervoice{00:31:20--00:32:10，老师用 Kasparov 对 Deep Blue 的评价做对照，强调现在看到的轨迹不再只是人工写好的搜索程序。讲义提醒：这是对行为形态的描述，不足以证明模型没有进行隐式搜索，也不证明轨迹忠实。}

\begin{importantbox}{“涌现”在这里应怎样理解}
涌现不是魔法标签，而是指训练目标没有逐条编码“先估规模、再找分解、最后验证”的脚本，模型却在规模、数据和训练后呈现这种可复用结构。研究问题应继续追问：哪些数据与奖励诱发它、在哪些分布失效、能否被 verifier 检查。
\end{importantbox}

\subsection{2025 算术题：长轨迹提供了什么}

题目要求用 1 到 10 各一次，只用加法和乘法得到 2025。简洁构造是
\[
(10\times4+5)\times(9\times3+8+7+2+1)=45\times45=2025.
\]
其中第一组构成 45，第二组也构成 45。slide 28 对比人工式简洁解和 2024 年 12 月 Gemini 2.0 thinking mode 的长轨迹：模型先注意 2025 是 $45^2$，再把目标转化为构造两个 45，并在很长的候选探索后给出答案。

\lecturefigure{slide-28.jpg}{2025 算术题：从规模判断到构造两个 45 的长推理轨迹}{28}

读图时先看右侧轨迹的\textbf{问题重表述}：大目标提示乘法，平方结构提示两个因子；再看搜索过程是否保持“每个数字只用一次”的约束；最后才看答案。长 CoT 的价值不是字数，而是把约束、启发式与中间目标显式化，使 verifier 可以逐步检查，也让后续 token 读取已经发现的 $45$ 结构。

\begin{warningbox}{长轨迹的三类幻觉}
第一，先猜到答案再编造合理过程；第二，中间违反约束但最终式子碰巧正确；第三，重复、绕圈或无效探索消耗大量 token。评测应同时检查最终表达式、数字使用集合、运算约束和轨迹成本。
\end{warningbox}

\subsection{学习发现过程，而不只是装入发现结果}

Rich Sutton 的引文强调：希望 agent 学会 discovery process，而不是只包含人类已经发现的答案。对 reasoning model，这意味着训练数据不能只有高质量终稿，还要提供可探索的问题分布、可靠反馈和从失败中恢复的机会。否则模型更像庞大的解答检索器，在新组合上仍缺少发现机制。

\lecturefigure{slide-29.jpg}{Bitter Lesson 视角：目标是学习发现过程}{29}

\teachervoice{00:32:10--00:39:00，老师用长算术轨迹说明模型会形成类似人类的启发式，但反复提醒 LLM 是 probabilistic model。课堂提示：可读步骤是工程资产，可用于验证与训练；“像人”只是一种描述，不应成为正确性的证据。}

搜索在这里并未消失，而是分成三层：参数中的模式压缩、token 序列中的局部探索，以及可选的外部 symbolic search/tool use。Q\&A 中老师进一步澄清，模型可以写 Python 或调用搜索器；这属于工具使用。研究“基础 reasoning ability”时可以单独考察无外部搜索的生成，但产品系统通常应把两者组合，而不是把 search 与 learning 设成意识形态对立。

\subsection{RL reasoning 的能力边界：不是所有任务都有自动 verifier}

slide 30 给出 RL finetuning 的核心优缺点：在答案可自动验证的任务上泛化较好，但大量开放任务没有唯一、便宜、可靠的评分函数。数学 exact match、代码测试、棋局输赢属于相对友好的区域；政策建议、长文研究、产品设计和社会判断往往有多目标、延迟反馈与价值冲突。

\lecturefigure{slide-30.jpg}{RL finetuning 的边界：自动可验证任务与开放任务}{30}

可把任务按 verifier 条件分成四类：
\begin{center}
\begin{tabularx}{0.96\textwidth}{p{0.19\textwidth}X X}
\toprule
类型 & 典型反馈 & 主要风险 \\
\midrule
唯一答案 & exact match、数值容差 & 解析错误、答案泄漏 \\
可执行结果 & 单元测试、模拟器回报 & 测试不完备、reward hacking \\
成对偏好 & 人类/模型 judge 选择 & judge 偏差、风格偏好 \\
长期开放目标 & 业务结果、社会影响 & 延迟归因、多目标冲突、不可逆成本 \\
\bottomrule
\end{tabularx}
\end{center}

\begin{knowledgebox}{Verifier frontier}
推理训练的下一阶段不是把已有 RL 配方机械移到所有领域，而是构造更可靠的过程反馈、可执行环境、反事实测试和多指标验收。没有这些基础设施，开放任务上的“奖励”很容易只测格式、讨好或短期代理。
\end{knowledgebox}

\subsection{为什么下一步是聚合与检索}

即使单条 reasoning path 已经很强，生成仍是随机过程：同一问题可能得到不同中间步骤和答案。slide 31 预告两个互补方向：aggregation 用多个候选估计答案概率，retrieval 从外部或上下文中补充相关知识。slide 32 再次提醒，LLM 是 next-token probability model，不是人类；“先生成轨迹再给答案”的自然流程仍可能在概率目标上错位。

\lecturefigure{slide-31.jpg}{进一步提升的两条路径：Aggregation 与 Retrieval}{31}

\lecturefigure{slide-32.jpg}{即使有 reasoning，next-token decoding 仍可能与最终答案目标不一致}{32}

\teachervoice{00:39:00--00:40:17，老师从 RL 的 verifier 限制转向 inference-time 改进：当单条轨迹不稳定时，不要假设一次生成足够；应利用多个样本的统计结构，或先找回相关知识再解题。}

\subsection{本章小结}

长 reasoning trace 能承载中间目标、约束和局部验证，但不能靠“像人”证明正确。RL 在自动可验证任务上具有优势，其扩展瓶颈是开放目标的 verifier。接下来要解决的，是单条随机轨迹与最终答案概率不一致的问题。

